\section{Cryptographic Design}\label{sec:crypto}

GhostVault defines no cryptographic algorithm of its own. It uses AES-256-GCM~\cite{nist-gcm},
HKDF-SHA256~\cite{rfc5869}, Argon2id~\cite{rfc9106} (or scrypt~\cite{rfc7914} where the
\code{cryptography} build has no Argon2id) and HMAC-SHA256~\cite{rfc2104}, all from the
\code{cryptography} library~\cite{pyca}, and SHA-256 for content checks.

\begin{figure*}[t]
\centering
% Fig.: the keys, what each is made from, and what each protects.
\begin{tikzpicture}[
  font=\footnotesize,
  box/.style={draw=black!65, fill=white, rounded corners=1.2pt, align=center, inner sep=2.4pt,
              minimum height=7mm, text width=#1},
  box/.default=30mm,
  dark/.style={box=#1, fill=black!78, draw=black!78, text=white},
  data/.style={box=#1, fill=black!5},
  arr/.style={-{Latex[length=1.4mm,width=1.1mm]}, draw=black!65, semithick},
  how/.style={font=\scriptsize, fill=white, inner sep=1pt, align=center},
  note/.style={font=\scriptsize\itshape, text=black!62, align=center,
              execute at begin node={\hyphenpenalty=10000\exhyphenpenalty=10000}},
  panel/.style={font=\scriptsize\scshape, text=black!60},
]
% ---- (a) one ghost --------------------------------------------------------------------------
\node[dark=78mm] (fk) at (0,0) {\textbf{file key}: 32 random bytes, one per ghost, kept in its record};
\node[data=34mm, below=12mm of fk.south west, anchor=north west] (chunks)
  {encrypted chunks $0,\dots,n{-}1$\\\scriptsize on GitHub};
\node[box=30mm, below=12mm of fk.south east, anchor=north east] (pk) {preview key};
\node[data=30mm, below=9mm of pk] (pv) {encrypted preview\\\scriptsize on GitHub and in the cache};
\draw[arr] (chunks.north |- fk.south) -- node[how] {AES-256-GCM} (chunks.north);
\draw[arr] (pk.north |- fk.south) -- node[how] {HKDF-SHA256} (pk.north);
\draw[arr] (pk) -- node[how] {AES-256-GCM} (pv);
\node[note, below=1.5mm of chunks, text width=40mm]
  {a fresh random 96-bit nonce per chunk, kept in the record; associated data: file id, index,
   count, size, padding};
\node[note, below=1.5mm of pv, text width=34mm] {padded: every preview is 32\,KiB\,+\,28\,bytes};
\node[panel, anchor=south west] at ([yshift=2mm]fk.north west) {(a) one ghost};

% ---- (b) recovery ---------------------------------------------------------------------------
\begin{scope}[xshift=90mm]
\node[dark=37mm] (pp) at (-20.5mm,0) {\textbf{passphrase}\\\scriptsize 12+ characters, or generated};
\node[dark=37mm] (rk) at (20.5mm,0) {\textbf{recovery key}\\\scriptsize 130 random bits, shown once};
\node[box=37mm, below=12mm of pp] (fkey) {recovery-file key};
\node[box=37mm, below=12mm of rk] (bkey) {backup key\\\scriptsize kept in \texttt{settings.json}};
\node[data=37mm, below=9mm of fkey] (file) {recovery file\\\scriptsize wherever the user keeps it};
\node[data=37mm, below=9mm of bkey] (bak) {repository backup\\\scriptsize \texttt{backup.gvrb} in each repository};
\draw[arr] (pp) -- node[how] {Argon2id, own salt} (fkey);
\draw[arr] (rk) -- node[how] {Argon2id, own salt} (bkey);
\draw[arr] (fkey) -- node[how] {AES-256-GCM} (file);
\draw[arr] (bkey) -- node[how] {AES-256-GCM} (bak);
\node[note, below=1.5mm of file.south east, anchor=north, xshift=2mm, text width=70mm]
  {both hold every complete record (file key, nonces, chunk map) and the repository list,
   never a token};
\node[panel, anchor=south west] at ([yshift=2mm]pp.north west) {(b) recovery};
\end{scope}
\end{tikzpicture}

\caption{The keys. (a) Each ghost has its own random key; its preview key is derived from it, so
chunks and preview never share a key. (b) A recovery bundle holds the records, with their keys,
encrypted under a key derived from a passphrase or a generated recovery key; each bundle key has
its own random salt. Black: secrets that are generated or chosen, never derived.}
\label{fig:keys}
\end{figure*}

\subsection{Keys}\label{sec:keys}
\figref{fig:keys} shows the keys and \tabref{tab:sizes} their sizes. A \emph{file key} is 32 bytes
from the operating system's random source (Python's \path{secrets.token_bytes}), one per ghost, stored in
its record. It encrypts that ghost's chunks and nothing else. The \emph{preview key} is
HKDF-SHA256 of the file key with the label \path{ghostvault/v6/subkey/preview}; it is derived when
needed and never stored. The keys of recovery bundles are derived with Argon2id from a passphrase
or a recovery key, each with its own random 16-byte salt (\secref{sec:recovery}).

Keys at rest are plain base64 in the records and settings, in files only the user can read. This is
deliberate. Wrapping them with a key kept in the operating system's keychain would protect nothing
against malware running as the user, which can read the keychain too, and losing that one keychain
entry would lose every ghost at once. The protection for a stolen computer is full-disk encryption;
backups of the data folder contain the keys and must be protected like it.

\subsection{Chunks}\label{sec:chunks}
A file is read in chunks of 512\,KiB of plaintext; a file of $s$ bytes has $\lceil s/524\,288
\rceil$ chunks, so an empty file has none. Each chunk is encrypted on its own with AES-256-GCM
under the file key and a fresh random 96-bit nonce; the ciphertext is the plaintext followed by a
16-byte tag. The nonce stays in the record; GitHub stores only the ciphertext.

Each chunk's associated data (\lstref{lst:aad}) names the file, the chunk's position, the number of
chunks, the file's exact size and whether the last chunk is padded. A chunk therefore cannot be
moved into another file, to another position, or into a file with more or fewer chunks without
failing authentication, and the size and padding flag in a record cannot be changed without every
chunk failing. The length used to strip the padding is thus authenticated, not taken on trust from
the record.

\begin{lstlisting}[basicstyle=\ttfamily\scriptsize,caption={Associated data of a chunk;
\code{p} or \code{u} says whether the last chunk is padded.},label={lst:aad}]
ghostvault/v6/<file id>/<index>/<count>/<size>/<p|u>

ghostvault/v6/5ada355b920792c5ba14be70737ffc5a/9/10/4860280/p
\end{lstlisting}

Nonces are random, and every encryption under a key uses a fresh one. With random 96-bit nonces,
NIST limits a key to $2^{32}$ encryptions~\cite[Sec.~8.3]{nist-gcm}. A file key is used once per
chunk of one file (a file of 1\,TiB has $2^{21}$ chunks), a preview key once or a few times, a
recovery-file key once, and the backup key once per backup version, so every key stays many orders of magnitude
below that limit.

Records written by earlier versions of GhostVault (record format 5) bind the file, the position and
the count, as \path{ghostvault/v4/<file id>/<index>/<count>}. They are still read, and restored
with that associated data; new records always use the form above.

\subsection{Padding}\label{sec:padding}
By default (\emph{Enhanced privacy} in the window), the last chunk's plaintext is filled with
random bytes up to the full 512\,KiB before it is encrypted, so every chunk of a ghost has the same
ciphertext size, 524\,304 bytes. No chunk is ever added: a file whose size is a multiple of 512\,KiB
is unchanged, and an empty file has no chunks. Restoring decrypts each chunk, checks its length,
and keeps exactly the recorded size. Each ghost records whether it was padded, and \emph{Storage
efficiency} turns padding off for new ghosts. \tabref{tab:padding} shows what GitHub stores.

\begin{table}[t]
\caption{Bytes of chunk ciphertext stored on GitHub}\label{tab:padding}
\centering\footnotesize
\begin{tabular}{@{}rrrr@{}}
\toprule
\textbf{File size (bytes)} & \textbf{Chunks} & \textbf{Padded} & \textbf{Not padded} \\
\midrule
0 & 0 & 0 & 0 \\
1 & 1 & 524\,304 & 17 \\
102\,400 & 1 & 524\,304 & 102\,416 \\
524\,288 & 1 & 524\,304 & 524\,304 \\
524\,289 & 2 & 1\,048\,608 & 524\,321 \\
5\,242\,881 & 11 & 5\,767\,344 & 5\,243\,057 \\
\bottomrule
\end{tabular}

\vspace{2pt}\raggedright\scriptsize Measured with \code{tools/measure.py}; each chunk adds a 16-byte
tag. Padded, every chunk is 524\,304 bytes.
\end{table}

Padding hides a file's exact size, not its approximate size: GitHub still sees the chunk count,
which gives the size to within 512\,KiB. The cost is at most 512\,KiB per ghost on GitHub.

\subsection{Checks on restore}\label{sec:checks}
A restore accepts a chunk only after five checks, all done before anything is written under the
file's final name:
\begin{enumerate}
\item the Git blob id is recomputed from the downloaded bytes and must equal the recorded id;
\item the SHA-256 of the ciphertext must equal the one recorded at upload;
\item the GCM tag must verify, with the associated data above;
\item for records of format 6, the plaintext length must be exactly what the recorded size and
padding imply;
\item after the last chunk, the whole file's SHA-256 must equal the original's.
\end{enumerate}
Git's ids are SHA-1 hashes. GhostVault uses them to name and find what GitHub stores and to detect
damage, not to resist an adversary: integrity against an adversary rests on the GCM tag and the
SHA-256 checks.

\subsection{Previews}\label{sec:previews}
For an image (JPEG, PNG, GIF, WebP or BMP, recognised by extension and decoded only in those
formats), GhostVault makes a preview of at most 200\,$\times$\,150 pixels, encoded as JPEG in at
most 32\,KiB\,$-$\,4 bytes. Decoding is bounded: images over 50 megapixels are refused from their
header, before any pixel is read, and JPEGs are decoded at reduced scale. The preview is encrypted
as soon as it exists and is never written in plaintext:
\begin{align*}
  \mathit{blob} = n \,\|\, \mathrm{AES\text{-}GCM}_{k_p}\bigl(&n,\; \ell \,\|\, \mathit{image} \,\|\,
  \mathit{pad},\\ &\text{``}\texttt{ghostvault/v6/preview/}\mathit{id}\text{''}\bigr),
\end{align*}
where $k_p$ is the preview key, $n$ a fresh 12-byte nonce, $\ell$ the image length in four bytes,
and the random padding brings the plaintext to 32\,KiB. Every preview blob therefore has the same
size, 32\,KiB plus 28 bytes. It is uploaded under its own random path to the repository of the
ghost's first chunk, committed with the chunks, and kept in a local cache in the same encrypted
form; it is decrypted in memory to be shown. Hiding a ghost's preview deletes only the local copy.
A preview reveals that a ghost has one, and nothing else.

\begin{table}[t]
\caption{Fixed sizes and parameters}\label{tab:sizes}
\centering\footnotesize
\begin{tabular}{@{}lr@{}}
\toprule
\textbf{What} & \textbf{Value} \\
\midrule
file key, preview key, bundle keys & 32 bytes \\
nonce, for every encryption & 12 bytes \\
authentication tag & 16 bytes \\
chunk plaintext & 524\,288 bytes \\
chunk ciphertext, padded & 524\,304 bytes \\
preview blob & 32\,796 bytes \\
file id & 16 random bytes (32 hex) \\
chunk path & \code{ghostvault/<2 hex>/<32 hex>.bin} \\
Argon2id: memory, passes, lanes & 256\,MiB, 3, 4 \\
scrypt fallback: $N$, $r$, $p$ & $2^{17}$, 8, 1 \\
bundle salt & 16 bytes \\
bundle padding unit & 64\,KiB \\
recovery key & 130 bits + 2 check characters \\
\bottomrule
\end{tabular}
\end{table}
